\section{Conservation and the diagonal-stress spectral bound}
\label{sec:physical-ward-stress-cone}

We now supply the conservation estimates required in
Section~\ref{sec:physical-ward-packet-prices}. The stress tensor $J_{ab}$
is real, symmetric, conserved, and covariant under signed coordinate
permutations. Its derivatives are the actual insertion derivatives of the
same hierarchy. The argument requires neither full rotation covariance nor
a value of a correlation at a collision.

Write $K_{abcd}(x)=S(J_{ab}(x),J_{cd}(0))$ for $x\ne0$.
Choose one positive integer $D$ large enough that these finitely many
kernels and their first derivative kernels are bounded by
$C\min(1,\|x\|_\infty)^{-D}$. Such an integer exists by the separated
polynomial derivative bounds. In this section a test functional is required
only on products $\prod_i h_i(x_i)$ of one-dimensional Schwartz functions.
We do not assert a collision extension with additional positivity properties.

\subsection{A symmetric extension and its contact terms}

Choose compactly supported smooth functions $g_r$ with
$g_r^{(k)}(0)=\delta_{rk}$ for $k<D+1$, and
$g_r(-x)=(-1)^r g_r(x)$. For example use $x^r/r!$ times a fixed even
cutoff equal to one near zero. Put
$\Pi_D h=\sum_{r<D}h^{(r)}(0)g_r$. Define
\[
 R_{abcd}(h)=\int K_{abcd}(x)
       \left(\prod_i h_i(x_i)-\prod_i\Pi_Dh_i(x_i)\right)dx.
\]
The difference of products is a finite telescoping sum, each of whose
terms contains one factor vanishing to order $D$ at zero. Taylor's theorem
bounds that factor by $C|x_i|^D\leq C\|x\|_\infty^D$ near zero.
Thus the integrand is bounded there. Away from zero the kernel is bounded
and each product is integrable. This proves absolute integrability and
multilinearity. If one $h_i$ vanishes to order $D$, its Taylor projection
is zero and $R(h)$ equals the ordinary integral of the kernel.

Use the same cutoff in all four coordinates. The product subtraction
commutes with signed coordinate permutations. The original kernel's tensor
law, symmetry in each index pair and exchange of the two pairs therefore
pass to $R$. Pair exchange uses translation, permutation of insertions and
inversion $x\mapsto-x$; the rank-four tensor has positive inversion sign.

Let $\Delta_{bcd}(h)=\sum_a R_{abcd}(h[a\leftarrow h_a'])$.
If one factor $h_i$ vanishes to order $D+1$, integration by parts is
legitimate and conservation gives $\Delta_{bcd}(h)=0$.
Indeed the original test and each differentiated test cancel the stated
collision powers; integrate first along a coordinate line whose transverse
coordinates are nonzero, where the kernel is smooth, and then use Fubini.
Those exceptional transverse coordinates have Lebesgue measure zero.
The Schwartz decay and the same integrable bounds justify the integration
by parts. This avoids imposing an unproved boundary value at the origin.

Replacing the four factors successively by their Taylor projections of
order $D+1$ now shows that $\Delta$ depends on finitely many jets at zero.
We normalize its linear jets. Let $h^{(a)}$ have factor $g_1$ in coordinate
$a$ and $g_0$ elsewhere, and set $c_{abcd}=\Delta_{bcd}(h^{(a)})$.
Reflection of a coordinate appearing an odd number of times among
$a,b,c,d$ changes this number's sign, so $c_{abcd}=0$ unless the indices
occur in pairs. In the three possible pair patterns, coordinate permutation
and symmetry in $c,d$ show
$c_{abcd}=c_{bacd}=c_{cdab}$. For instance, the patterns $aabb$, $abab$
and $abba$ reduce respectively to swapping the two coordinate names,
swapping them followed by the last-pair symmetry, and those same two
operations. Thus the correction
\[
 \widetilde R_{abcd}(h)=R_{abcd}(h)-c_{abcd}\prod_i h_i(0)
\]
retains all the stated tensor symmetries and all collision-free evaluations.
Its divergence has zero linear jets. Inversion makes that divergence odd,
so its jets of even total order are also zero. Its finite jet expansion
therefore contains only total orders at least three.

For real frequency profiles $f_i$ let
$T_{abcd}(f)=\widetilde R_{abcd}((\widehat f_i)_i)$, and write $P_iT$
for insertion of the multiplier $p_i$ in that frequency profile. Our
Fourier convention is $(\widehat f)'=-i\widehat{pf}$. Consequently
\[
 Q_{bcd}=\sum_a P_aT_{abcd}
\]
is $i$ times the divergence functional. If all profile widths are
$e^{w_in}$ with $w_i\leq a\leq0$, its finite jet expansion has exponent
$3a$. Indeed a jet of multiorder $\alpha$ is a fixed coefficient times
$\exp(n\sum_iw_i\alpha_i)$, and its total order is at least three.
Multiplication by one extra frequency coordinate adds at most $a$ to this
exponent. These are actual estimates for the normalized contact terms.

\subsection{Fixing one frequency band}

We need a second estimate: if the profile in axis $i$ is fixed, compactly
supported and vanishes near zero, and the other widths are $e^{w_jn}$,
$w_j\leq0$, then $T_{abcd}$ has every positive exponent $\delta$, provided
the other profiles have the reflection parities of these four tensor indices.
We give the details that make this estimate independent of a rotation law.

First extend the positive subtraction argument to an odd temporal profile
$f$. Define its correction by
\[
 B^{\rm odd}_{m,s}f(p)=p^{-1}B_{m,s}(pf)(p).
\]
The numerator vanishes near zero, so this is a smooth compactly supported
function. The residual has all moments below $2m$ equal to zero: its even
moments vanish by oddness, and each odd moment is an even moment of the
ordinary residual of $pf$. For $E>0$ its signed Laplace response is
\[
 \mathcal L_{-1}(f-B^{\rm odd}_{m,s}f)(E)
       =-2i\int\frac{R_{m,s}(pf)(p)}{p^2+E^2}\,dp.
\]
The positive-kernel formula still applies to the integral on the right.
For fixed odd $f$, $|f(p)|\leq L|p|$ by the fundamental theorem of
calculus. If its support is in $[-R,R]$, $R\geq1$, then
$|p f_w(p)|\leq Lp^2/w^2$ on $|p|\leq Rw$.
Every kernel factor is bounded by
$C_{m,s}(1+E)^{-2m}$ in product. When $E\leq w$, use
$p^2/(p^2+E^2)\leq1$; when $E\geq w$, use
$p^2\leq R^2w^2$ and $(p^2+E^2)^{-1}\leq E^{-2}$.
Integration over an interval of length $2Rw$ gives in both cases
\[
 |\mathcal L_{-1}(f_w-B^{\rm odd}_{m,s}f_w)(E)|
 \leq C(E+w)^{-1}(1+E)^{-2m}.
\]
The positive lower bound in \eqref{eq:packet-response-floor} therefore
majorizes this by $C'M_{m,s}\chi_w(E)$, uniformly for $0<w\leq1$ when
$s$ contains the supports. The assertion at $E=0$ follows by continuity.
For an even profile, choose an even nonnegative compactly supported
majorant $g\geq|f|$ and use positivity on $g+f$ and $g-f$ instead.
Thus either parity has a positive response majorant.

For two clusters $A,B$, the common projection-valued measure gives a
cross-measure estimate. If $|q|^2\leq uv$ with $u,v\geq0$, then
\begin{equation}\label{eq:stress-cross-price}
 |\mu_i^{B,A}(q)|\leq t\int u\,d\mu_i^B
                       +t^{-1}\int v\,d\mu_i^A,\qquad t>0.
\end{equation}
To prove it, first use a common admissible positive-time placement.
Relative to the sum of its two finite diagonal spectral measures, the
$2\times2$ matrix of densities is positive almost everywhere: test its
quadratic form on the countable dense set of vectors with rational real
and imaginary parts. Its off-diagonal entry is therefore bounded by the
geometric mean of the diagonal entries. Integrate, apply Cauchy--Schwarz
and then $\sqrt{UV}\leq tU+t^{-1}V$. Exponential undamping multiplies all
entries by the same positive function, so preserves the inequality.
Truncation and absolute integrability justify its use for the responses
here, including energy-zero atoms.

Suppose each cluster has both reflection signs among the four axes and
is on every coordinate hyperplane. In a chosen coordinate subtract enough
moments from its frequency profile, using the even or odd correction
according to the product of the two reflection signs. The previous
majorant, the mixed-sign estimate of
Section~\ref{sec:physical-ward-packet-prices} for each diagonal measure,
and \eqref{eq:stress-cross-price} give exponent $\delta$ for the tested
cross correlation. All transverse profiles can first be majorized by fixed
even nonnegative profiles. There are only two diagonal measures and finitely
many polarization integrability conditions, so choose the subtraction order
after all of them. This proves the estimate with one subtracted coordinate.

Now keep a profile in coordinate $i$ fixed and away from zero. At another
coordinate $j$ split the shrinking profile into its moment-zero residual
and its finite correction. The residual is controlled by the preceding
estimate. Each correction is a finite linear combination of fixed annular
profiles, with coefficients of exponent at most zero, since $w_j\leq0$.
For odd parity the additional division by $p_j$ is harmless on those fixed
annuli. Apply induction to the smaller set of shrinking coordinates in
each correction term. At the base every profile is fixed, so the functional
is constant. The correction order may depend on all profiles already fixed
at that stage; only finitely many such choices occur. No uniform choice
over an infinite family of profiles is required.

To apply this to stress components, an off-diagonal component $J_{ab}$,
$a\ne b$, already has both signs: it is odd in axes $a,b$ and even in
the other two. A diagonal component is replaced by $D_iJ_{aa}$, which
is odd in $i$ and even in the other axes. Move that derivative to the
fixed $i$-profile by integration by parts, dividing that frequency profile
by $p_i$. This is smooth because the profile vanishes near zero.
The first and second insertions give respectively the factors $-i$ and
$i$; doing both replacements composes these identities. Sufficiently many
moments in the subtracted $j$-profile justify the ordinary integrals and
integration by parts, as in the extension construction above.

Choose representative insertion tokens for these field labels. Average
over the four commuting reflection involutions with the desired signs.
This finite projection leaves the represented field unchanged and produces
actual clusters with those exact signs; their positions are still zero.
The two modified stress labels thus satisfy the mixed-sign hypotheses.
The preceding finite induction proves the asserted fixed-band estimate
for the original extended stress functional, with the required transverse
parities. It also applies after any fixed finite number of frequency
multiplications, with their corresponding extra scale factors.

\subsection{Three responses linked by conservation}

Fix a spectral axis $i$. All functionals in this subsection refer to the
same normalized extension. Define
\begin{align*}
 T_0&=T_{iiii},\\
 T_1&=\sum_{a,j\ne i}P_jP_aT_{aiji},\\
 T_2&=\sum_{a,b,c,j\ne i}P_jP_aP_bP_cT_{cj ab}.
\end{align*}
The indices in the last expression are grouped as the pairs $(c,j)$ and
$(a,b)$. Symmetry in and between the pairs, and
$Q_{bcd}=\sum_aP_aT_{abcd}$, give
\begin{equation}\label{eq:stress-conservation-chain}
 T_1=P_i^2T_0+C_1,\qquad T_2=P_i^2T_1+C_2.
\end{equation}
More explicitly,
\[
 C_1=\sum_{j\ne i}P_jQ_{iji}-P_iQ_{iii},\qquad
 C_2=\sum_{j,a,b\ne i}P_jP_aP_bQ_{jab}
                    -P_i\sum_{j,a\ne i}P_jP_aQ_{aji}.
\]
Expand the two divergence sums into their $i$ term and their remaining
terms to verify these identities; commuting the frequency multipliers and
exchanging the tensor pairs identifies the displayed $T_1,T_2$.
When all four widths are at most $e^{an}$, the contact bounds above give
exponents $4a$ and $6a$ for $C_1,C_2$, respectively.
They also show that $C_1$ vanishes if the $i$-profile has zero moments
through order $D+1$, and $C_2$ if it has zero moments through order $D+3$:
multiplication shifts the required moment order by at most one or three.

Fix $\delta>0$, $a\leq0$, and three even, nonnegative, smooth compactly
supported transverse profiles $f_j$, $j\ne i$. Keep their shapes fixed,
with width exponents $w_j\leq a$. The fixed-band result, with the $i$-profile fixed away
from zero, gives exponents
\[
 \delta,\qquad\delta+2a,\qquad\delta+4a
\]
for $T_0,T_1,T_2$: their respective sums contain zero, two and four
transverse frequency multipliers. The parity hypotheses hold because each
multiplier flips the sign in its coordinate, cancelling the corresponding
stress index sign.

Let $w=e^{an}$ and $\varphi_k(p)=p^{-2k}\varphi(p)$ as in the previous
section. Write $h_{p,w,k}=(\varphi_k)_{pw}$, for $p=1,2$, and define
\[
 A_{p,k}(n)=w^{-2k}\operatorname{Re}
       T_k(f_1^{(n)},\ldots,h_{p,w,k},\ldots,f_4^{(n)}),
       \qquad k=0,1,2.
\]
The temporal slot is $i$. Since $p_i^2h_{p,w,k+1}=(pw)^2h_{p,w,k}$,
\eqref{eq:stress-conservation-chain} implies that both
\[
 A_{p,1}-p^2A_{p,0},\qquad A_{p,2}-p^2A_{p,1}
\]
have exponent $2a$, hence exponent $x=\delta+2a$.
Here the normalized first defect costs $w^{-2}e^{4an}$, and the second
costs $w^{-4}e^{6an}$, making the exponent explicit.

Choose one subtraction order $m$ greater than all the finitely many moment
orders just used. Replace $h_{1,w,k}-h_{2,w,k}$ by its moment-subtracted
difference at scale $8$, and call its normalized real response $F_k$.
The correction coefficients are differences of the moments at widths $w$
and $2w$. Their constant moments cancel. Each remaining even moment has
order at least two, so its coefficient has exponent at most $2a$.
There are finitely many fixed correction profiles, all away from zero.
Their fixed-band costs are $\delta+2ka$ for $T_k$, and the normalization
$w^{-2k}$ removes the latter $2ka$. It follows that
\[
 F_k-(A_{1,k}-A_{2,k})=O(e^{xn}),\qquad k=0,1,2.
\]
No subtraction order is sent to infinity.

The three signs are essential. Take the actual reflection-even cluster
representing $J_{ii}$ and its positive spectral measure. For a sufficiently
moment-zero temporal profile, the extension agrees with the ordinary
spectral integral. The contact terms in
\eqref{eq:stress-conservation-chain} vanish there; multiplication by
$p_i^2$ inside the profile becomes multiplication by $-E^2$ in its Laplace
response, by twice integrating by parts (the zero moments remove the
boundary terms). Thus, writing $\Gamma_{m,w,k}(E)$ for the nonnegative
dyadic response from the previous section,
\[
 F_k=\tau w^{-2k}\int(-E^2)^k\Gamma_{m,w,k}(E)
             \prod_{j\ne i}(f_j)_{e^{w_jn}}(p_j)\,d\mu_i^{J_{ii}}(p),
\]
where $\tau>0$ is the fixed Fourier normalization. Hence $F_0\geq0$,
$F_1\leq0$ and $F_2\geq0$.
The four conservation defects and these three correction defects are
exactly the seven errors in the algebraic estimate of
Section~\ref{sec:physical-ward-packet-prices}. It gives
$F_0=O(e^{(\delta+2a)n})$. The positive dyadic kernel lower bound and
multiplication by $w^2$ now prove \eqref{eq:packet-cone-required}.
In particular the choice $a=-99$, $\delta=1$ gives exponent $-395$ for
the actual diagonal-stress measure, as required by the exterior packet.
